% =====================================================================
%  Section 1 -- Introduction
% =====================================================================

\section{Introduction}
\label{sec:intro}

Floods are the most frequent and economically damaging natural disasters globally, accounting for 38.7\% of natural disasters and affecting 43\% of disaster-affected populations between 2000 and 2009~\cite{du2010health}. Rapid urbanisation continues to amplify exposure~\cite{hemmati2020urban}. Because traditional hydrological and hydraulic models suffer from high computational demand, coarse spatial resolution, and limited real-time adaptability~\cite{kumar2023flood}, deep-learning segmentation has emerged as a leading approach to automated flood-extent delineation from satellite, drone, and ground imagery~\cite{bentivoglio2022deep}.

Recent benchmark studies in this space have converged on encoder--decoder architectures --- UNet~\cite{ronneberger2015unet}, SegNet~\cite{badrinarayanan2017segnet}, and DeepLabV3$+$~\cite{chen2018encoder} --- paired with strong feature-extraction backbones, most commonly EfficientNet~\cite{tan2019efficientnet} and ResNet~\cite{he2016deep}. Karc{\i} \emph{et al.}~\cite{karci2026deep} provide the most comprehensive systematic comparison of these combinations on the public Flood Semantic Segmentation Dataset (FSSD), reporting that UNet $+$ EfficientNet achieves the highest segmentation performance (IoU 0.927, F1 0.957, accuracy 0.971) and conclude that this combination ``is ideal for near real-time operational deployment.''

A careful reading of the broader flood-segmentation literature, however, reveals that this operational-deployment claim has not actually been evaluated. Across the 26 related works surveyed by Karc{\i} \emph{et al.}~\cite{karci2026deep} and adjacent studies on UAV~\cite{hernandez2022flood}, SAR~\cite{pechmay2023flooded}, and multi-modal flood mapping~\cite{chamatidis2024vision}, every paper reports segmentation accuracy on academic benchmarks; \emph{none reports model parameters, FLOPs, per-platform inference latency, or post-quantization accuracy.} This omission is consequential. UNet $+$ EfficientNet-B7 --- the configuration favoured in~\cite{karci2026deep} --- has approximately 63\,M parameters and 50\,GFLOPs at $256{\times}256$ input resolution. Such a model cannot run on the drones, embedded sensors, Raspberry-Pi-class community stations, or offline tablet apps that humanitarian organisations such as UN-SPIDER and the ICRC actually deploy in flood zones, where reliable internet connectivity is the exception rather than the rule~\cite{munawar2021review}.

\subsection{Contributions}
\label{sec:contributions}

This paper makes three contributions.

\begin{enumerate}
\item \textbf{First systematic multi-platform edge-deployment benchmark for flood segmentation.} We report FP32 and INT8 inference latency for three lightweight UNet students on two deployment targets --- Apple M2~Pro via ONNX Runtime CPU, and browser WebAssembly SIMD via ONNX Runtime Web 1.17 --- together with per-model parameter counts, FLOPs at $256{\times}256$, ONNX disk size, and FP32-to-INT8 size-reduction factors. ONNX serves as the deployment hub: every measured latency reflects the same exported artifact running under a different execution provider. We benchmark these two commodity software runtimes on a single Apple~M2~Pro host; measurement on physically constrained edge devices (Raspberry~Pi, mobile SoCs, embedded ARM) is out of scope for this study and deferred to future work (Section~\ref{sec:limitations}).

\item \textbf{Negative empirical finding on knowledge distillation for FSSD-class flood segmentation.} Across three architectures $\times$ four KD configurations $\times$ three folds, no KD configuration improves on the no-KD baseline. ImageNet-pretrained lightweight students match or exceed the UNet$+$EfficientNet-B0 teacher without distillation, and combined response $+$ feature KD strictly underperforms the no-KD baseline on every student. We provide a mechanistic explanation (saturated 663-image benchmark; strong ImageNet priors leave no headroom for the teacher's soft targets). The effect is directionally consistent --- across all three architectures no KD configuration attains a higher mean IoU than its no-KD baseline, and the best student exceeds the teacher in every one of the three folds --- but the fold-level design ($n=3$) is underpowered for a conventional significance test (minimum reportable two-sided Wilcoxon $p=0.25$), so we report KD as non-beneficial on this saturated benchmark rather than claiming a statistically powered effect (Section~\ref{sec:setup-stats}).

\item \textbf{Architecture-dependent post-training quantization stability characterisation.} Under an identical ONNX static QDQ INT8 recipe (Percentile 99.999\% calibration, per-tensor activations, per-channel weights restricted to Conv/Gemm/MatMul), EfficientNet-Lite0 quantizes cleanly ($\Delta \mathrm{IoU} = -0.043 \pm 0.016$) while MobileNetV3-Small and MobileViT-XXS suffer 0.45--0.46 IoU collapses with cross-fold standard deviations of 0.19--0.34. We trace the failure modes to (i)~the HardSwish activations in MobileNetV3-Small that lack bounded calibration ranges, and (ii)~the LayerNorm and self-attention blocks in MobileViT-XXS that produce rank-deficient calibration statistics under per-channel quantization. The implication for practitioners: \emph{bounded-activation backbones (ReLU6) post-training-quantize reliably for flood segmentation; HardSwish and attention-based backbones require quantization-aware training (QAT).}
\end{enumerate}

We make our trained checkpoints, the ONNX FP32 and INT8 exports, the per-platform benchmarking harnesses, and a Zenodo-archived release of the full code repository publicly available.

\subsection{Scope and pivot from initial framing}
\label{sec:scope-pivot}

This work was originally framed as a knowledge-distillation recipe for compressing a heavy flood-segmentation teacher. The empirical results did not support that framing, and we report the negative-KD finding rather than search for a configuration that would justify the original framing. The contribution pivots to the per-platform edge benchmark and the architecture-dependent PTQ characterisation. We hope this orientation is more useful to practitioners than another paper-grade KD recipe, and that the negative result helps future researchers avoid investing compute in distillation on saturated benchmarks where strong ImageNet pretraining has already closed the gap.

The remainder of the paper is organised as follows. Section~\ref{sec:related} reviews relevant prior work on flood segmentation, lightweight segmentation backbones, knowledge distillation, and post-training quantization. Section~\ref{sec:methods} details the FloodLite methodology. Section~\ref{sec:experimental-setup} describes the experimental setup, including the 3-fold protocol and per-platform latency rigging. Section~\ref{sec:results} reports results across Tables~\ref{tab:t1-footprints}--\ref{tab:t5-latency} and Figures~\ref{fig:f2-pareto}--\ref{fig:f4-throughput}. Section~\ref{sec:discussion} discusses the negative-KD finding, per-platform deployment recommendations, and the architecture-dependent PTQ insight. Section~\ref{sec:limitations} lists limitations. Section~\ref{sec:conclusion} concludes.
